\documentclass[../../../../SoftwareRequirementsSpecification.tex]{subfiles}
\begin{document}
% Richiede le macro definite in src/macros/useCaseMacros.tex
% (incluse nel preambolo di SoftwareRequirementsSpecification.tex)

\ucstart
\begin{xltabular}{\textwidth}{|l|X|}
  \hline
  \ucid{PT-01} \\ \hline
  \ucname{Registrazione PT} \\ \hline
  \ucactor{PT} \\ \hline
  \ucdescription{Il PT crea da solo il proprio account per usare
    l'applicativo.} \\ \hline
  \ucprecond{L'utente non ha un account nel sistema.} \\ \hline
  \ucpostcond{Il PT ha un account nell'applicativo. Può accedere
    tramite il caso d'uso U-01. Non ha ancora nessun atleta
    associato.} \\ \hline
  \ucmainflow{
    \item L'utente apre la pagina di registrazione.
    \item Il sistema mostra il form di inserimento dati (nome,
      cognome, email e password).
    \item L'utente inserisce i dati.
    \item L'utente preme il pulsante "Registrati".
    \item Il sistema verifica che i dati siano validi e che l'email
      non sia già presente nel database.
    \item Il sistema crea l'account del PT, memorizzando l'impronta
      della password e non la password stessa.
    \item Il sistema mostra una pagina di conferma e indirizza
      l'utente al caso d'uso U-01 (Login).
  } \\ \hline
  \ucaltflow{
    \textbf{5a. Email già registrata:} \newline
    - Il sistema mostra un messaggio di errore ("Email già
    registrata"). \newline
    - Il flusso riprende dal passo 2. \newline
    \textbf{5b. Dati non validi:} \newline
    - Il sistema mostra un messaggio di errore ("Controllare i dati
    inseriti") e indica i campi non validi (email in formato non
    valido, nome o cognome vuoti, password che non rispetta i
    requisiti minimi di lunghezza e composizione). \newline
    - Il flusso riprende dal passo 2.
  } \\ \hline
  \ucnote{
    Il PT sceglie la propria password al momento della
    registrazione. Il sistema non la memorizza in chiaro, ma solo la
    sua impronta calcolata con una funzione di hash. \newline
    Al termine della registrazione il sistema non autentica l'utente,
    ma lo indirizza al login: le credenziali appena create vanno
    usate tramite il caso d'uso U-01. \newline
    Il messaggio di errore del flusso 5a è lo stesso in due casi:
    quando l'email appartiene già a un altro PT, e quando appartiene
    a un atleta già registrato dal proprio PT. Il campo
    \texttt{email} di \texttt{User} è infatti univoco su tutta la
    gerarchia degli utenti. Per questo motivo, uno stesso utente non
    può avere sia il ruolo di PT sia quello di Atleta.
  } \\ \hline
\end{xltabular}
\end{document}
